% ! TeX root = ../erc-B1.tex
\begin{refsection}

% Add \subsection's as needed. It is possible to refer to a part of the
% B2 using \input{../part-B2/parts/file}; within this file, \ifbone can
% be used to determine if one is in the B1 or in the B2

\todo{
    \textbf{Part I of the Scientific Proposal } should present the envisaged research and it should:
    \begin{itemize}
        \item \textit{lay out the current state of knowledge, }
        \item \textit{explain the scientific question and the objectives of the project, and}
        \item \textit{present the overall approach or research strategy to reach the goals of the
        project.}
    \end{itemize}
    %
    Part I should convince the evaluation Panel that it presents an original
    and creative idea addressing an important question in the respective
    research field(s). Furthermore, it should substantiate how the
    project will advance the frontier of knowledge, and what contribution
    it will make to the research field(s) i.e. what may be changed,
    opened, challenged or how the results of the work will alter the
    current understanding of the field.
    %
    \textbf{At Step 1, only Part I and the Curriculum Vitae (CV) and Track Record
    (see below) is assessed by the evaluation Panel.}
    %
    It forms the basis for the Panel’s decision whether it chooses to
    evaluate the proposal in the next step. Therefore, all essential
    information must be covered in this section.
}

\todo{
    References
    to literature should also be included. Please use a reference style
    that is commonly used in your discipline such as American Chemical
    Society (ACS) style, American Medical Association (AMA) style, Modern
    Language Association (MLA) style, etc. \underline{and}
    that allows the evaluators to easily retrieve each reference.
}

% Add \nocite{} for any citation from the CV or EATR to add to the
% bibliography

\pagebreak
\defbibheading{subsection}[\bibname]{\subsection*{#1}}
\printbibliography[prenote=bolditalics,heading=subsection]

\end{refsection}
